\section{Conclusions}\label{sec:conclusion}
We have examined how idealized foreground cuts affect the H\,I-weighted genus and its redshift-dependent response to AP mappings. Four Jiutian 1G snapshots at $0.10\le z\le0.59$ provide H\,I-weighted fields smoothed at $20\,h^{-1}\mathrm{Mpc}$. We compare their genus curves, subtract an intrinsic-evolution template and examine horizontal and wedge cuts across eleven slopes.

Foreground removal changes the topology substantially: the horizontal cut increases the fitted genus amplitude by about 4.6 while suppressing the density variance. Evolution across the sampled redshifts is modest, but endpoint changes alone omit intermediate variation. Incorrect mappings leave redshift-dependent residuals after template subtraction. Parameter-region areas, however, change under mesh and mask choices and parameter-space smoothing, so the current area ratios and sampled wedge minimum are not robust measures of cosmological precision.

These results motivate H\,I topology as a complement to conventional clustering statistics in the presence of foreground mode loss. Applying this approach to the SKA and related surveys will require realistic mock observations and a stable measurement definition \citep{Wang2015HITopologySKA}; no survey precision forecast is inferred here.

The present cuts omit instrumental response and residual foreground contamination. The intrinsic template depends on the simulation and H\,I population model, and the single-simulation covariance lacks independent calibration. The existing numerical comparisons expose sensitivity but do not demonstrate convergence or select a replacement estimator. Future work should establish numerical stability, propagate template uncertainty and test realistic foreground and survey treatments before quantitative cosmological inference.

\section*{Acknowledgements}
This work was supported by the National Natural Science Foundation of China (grant No. 12373005). We also wish to acknowledge the Beijing Super Cloud Center (BSCC) and Beijing Beilong Super Cloud Computing Co., Ltd (\url{http://www.blsc.cn/}) for providing HPC resources that have significantly contributed to the research results presented in this paper.

\section*{Data availability}
\todo{Confirm access conditions for the simulation catalogue and the repository or contact route for derived data and analysis code.}
